\section{Detailed measured results and evidence audit}
\label{m:results}

\textbf{Status of this section.} Capability, reading-comprehension, development-decoder, and long-context confirmation panels are measured. The decoder was frozen before confirmation; the complete paired comparison contains 105,600 validated outcomes across both models.

\subsection{Identity of the evaluated checkpoint}
\label{m:identity}

The evaluated payload is a local copy of the F2 step-14,000 release. Reading its state
dictionary directly, rather than quoting a model card, gives: \textbf{1{,}952 stored
tensors}, all in \texttt{float32}, totalling \textbf{4{,}091{,}581{,}440 stored
elements} ($4.0916\times10^9$). This independently reproduces the $\sim$4{,}091.6M element
count the protocol left to be obtained from the payload, and it is the \emph{stored} count:
the instantiated bf16 model executes the same number, and we report the element count
rather than a "trainable parameter" count that could differ by buffers.

The tensor inventory also confirms two structural claims that the protocol asks to be
checked rather than assumed. There are \textbf{zero} keys under \texttt{loop.} and
\textbf{zero} keys naming a latent, FiLM, or conditioning module, so the evaluated object
is the plain bidirectional denoiser with no depth recycling and no condition interface ---
it is not the later looped model, and no random module has been attached. The two mixer
directions are stored separately and equally: \texttt{attn\_fwd} and \texttt{attn\_bwd}
hold $0.9340$B elements each, $1.8681$B together, so a forward-only evaluation of the same
geometry would execute $3.1575$B elements and the reverse stream is $22.8\%$ of the
payload. Embedding and output projection are $167{,}772{,}160$ elements each
($65536\times2560$).

The local FP32 payload is pinned by \texttt{sha256}
\texttt{cb45226f\ldots f9bf19}, $16{,}367{,}166{,}427$ bytes, at step 14{,}000 with
$14{,}680{,}064{,}000$ tokens seen in its metadata. The public release is pinned to
revision \path{9d2b9042d7ab32e5d706479373a5437b1a29fb4b}. Its BF16 export has
$8{,}183{,}369{,}960$ bytes and SHA-256
\path{0e98d75dab14b16feffea4a6b692550cdd558b676d5058bcde0495b5b838d231}.
Reconstructing that export from the local tensors reproduces this digest exactly.
Thus the local FP32 storage and public BF16 export are intentional representations,
and the evaluation-time BF16 conversion is identical to the released tensor payload;
public weight identity is no longer an unresolved qualification.

\paragraph{Verifiable export audit.}
The audit fetched only the original $207{,}080$ header bytes, including the eight-byte
length prefix, from the revision-pinned safetensors file. It checked all $1{,}952$
keys and shapes against the local CPU-memory-mapped checkpoint, checked contiguous
nonoverlapping data offsets and total file size, and streamed SHA-256 over the original
header followed by locally converted tensors in the header's declared dtype and offset
order. The resulting digest equals the public revision's LFS SHA-256; no public weight
payload download or temporary 8\,GB export was required. The source package includes
\path{lrwkv_evidence/verify_f2_export.py} and
\path{results/f2_export_audit/audit.json} with the header digest and verification receipt.
This establishes export correspondence, not source-code or sampler equivalence.

No historical validation split or group-size decoder was recovered; the grouped
$g{=}4$ sampler named by the card remains a reproduction target. The capability
readings below use causal span likelihood without executing a sampling decoder.

The GPU runtime uses the FLA RWKV implementation, which emits a warning recommending comparison with the official RWKV implementation. The present qualification covers artifact identity, tokenizer IDs, target masking, and repeated-call/cache parity; it does not establish independent official-kernel equivalence. Generated behavior and timings are therefore specific to the recorded runtime.

\subsection{Capability: full MMLU}
\label{m:mmlu}

Table~\ref{tab:mmlu} reports the complete MMLU test panel, 14{,}042 items, under a one-token label convention whose prompt tree and label encoding are frozen artifacts on disk: the available RWKV and GLA tree manifests carry $\texttt{template\_sha256}=$\texttt{84ea26ef\ldots 213dc}, which is the SHA-256 of the prompt-template literal in the panel's own builder module, together with $\texttt{convention}=\texttt{one\_token\_answer\_label}$ and $\texttt{choice\_start}=L-1$. Two scoring notes. The diffusion rows are scored under the arm name \texttt{fwdce} and the autoregressive rows under \texttt{raw}; these are two \emph{names for one statistic}, not two statistics --- both reach the same length-normalized span likelihood through the same shift, and the per-class assignment exists only because the wrong name either crashes or silently returns a different quantity. The rows are therefore comparable across arms. Separately, the run records carry \texttt{unpinned} in their registry, profile and condition-profile hash fields. That is a provenance gap about the \emph{harness revision} those runs were produced under; it is not a statement about the prompt or the scoring rule, which are pinned where the tree manifests above are available. Mamba-1, Qwen, and Llama tree manifests are absent in the current evidence archive: their exact unique-document coverage and shard/merged agreement are verified, but the original builder/template pin remains incomplete (\path{results/additional_baseline_audit.json}). We record it as what it is rather than as a caveat on the statistic. The prompt trees are tokenizer-specific, so a row is scorable
only against the tree built with its own tokenizer; the released RWKV-7 and diffusion rows
share the RWKV-world vocabulary, and the rows are labelled accordingly. No subset is
reported as the benchmark: a card's 200-question observation is not full MMLU, and neither
is a capped prefix.

\begin{table}[t]
\centering\small
\caption{Full MMLU test panel, 14{,}042 items, one-token label convention, causal span
likelihood. $n=14{,}042$ for every row; intervals are the recorded 95\% percentile bootstrap intervals over document records and \emph{exclude} execution variance (Section~\ref{m:variance}). The diffusion
rows are the evaluated F2 payload and its looped siblings; the released rows are fixed
public checkpoints of different size and pretraining exposure. Sizes are stored-element
counts where measured, otherwise model labels.}
\begin{tabular}{@{}llrr@{}}
\toprule
ID & Model & Accuracy & 95\% CI\\
\midrule
F2 & BiRWKV diffusion, step 14{,}000 (this work's object) & 0.469 & $[0.461,0.477]$\\
--- & BiRWKV diffusion, step 13{,}000 & 0.427 & $[0.418,0.435]$\\
--- & Looped variant, step 4{,}750 & 0.475 & $[0.466,0.483]$\\
--- & Looped variant, step 9{,}500 & 0.479 & $[0.470,0.487]$\\
--- & Looped variant, step 12{,}500 & 0.467 & $[0.459,0.475]$\\
--- & Looped variant, step 15{,}000 & 0.451 & $[0.442,0.459]$\\
--- & Looped variant, step 18{,}500 & 0.458 & $[0.449,0.465]$\\
--- & Looped variant, step 19{,}500 & 0.462 & $[0.454,0.470]$\\
--- & Looped variant, step 20{,}000 & 0.453 & $[0.445,0.460]$\\
--- & Looped variant, step 20{,}500 & 0.452 & $[0.444,0.460]$\\
\midrule
R0 & RWKV-7 World3 2.9B (released) & 0.546 & $[0.538,0.554]$\\
--- & RWKV-7 World3 1.5B (released) & 0.420 & $[0.412,0.428]$\\
--- & RWKV-7 0.4B (released) & 0.260 & $[0.253,0.268]$\\
L0 & GLA-1.3B-100B (released) & 0.230 & $[0.223,0.237]$\\
--- & Mamba-1 2.8B (released) & 0.258 & $[0.251,0.265]$\\
--- & Qwen2.5-3B (released) & 0.653 & $[0.645,0.661]$\\
--- & Llama-3.2-3B (released) & 0.541 & $[0.534,0.549]$\\
\bottomrule
\end{tabular}
\label{tab:mmlu}
\end{table}

Two observations follow, and both are of the "practical comparison of fixed releases"
kind the design section warns against over-reading. The evaluated F2 payload reaches
$0.469$ on full MMLU, which is $7.7$ points below the released causal RWKV-7 of comparable
size ($0.546$) and $7.2$ points below Llama-3.2-3B, while exceeding the released
RWKV-7 1.5B ($0.420$) and the released 0.4B model ($0.260$). The Mamba-1 2.8B released checkpoint reaches $0.258$, which is $21$ points below
F2; note that the comparable Mamba-2 release named in Table~\ref{m:arms} has no reading
here at all, and the two architectures are not interchangeable. These are descriptive differences between releases differing in
pretraining data, training objective, tokenizer and exposure, so they bound the practical
comparison and isolate nothing.

Within the diffusion lineage the capability panel is \emph{flat and non-monotone} in
training step: the best full-MMLU reading in the ladder is step 9{,}500
($0.479$), and later steps are between $0.451$ and $0.467$. Training longer does not
improve this panel, and the spread across steps ($2.8$ points) is the same order as the
difference between the early and late ladder. A capability panel that does not move with
training step cannot be the evidence for a method whose claimed benefit is long-context
retrieval; it can only bound the cost of the architecture, which is what we use it for.

\subsection{Reading comprehension: full HellaSwag}
\label{m:hellaswag}

Table~\ref{tab:hellaswag} reports the complete 10{,}042-item HellaSwag validation panel.
The pre-existing readings for this benchmark in our earlier records were produced with a
1{,}000-item cap, and because the item files are task-grouped that cap returned one
activity group rather than a random tenth; one arm reads $0.573$ at $n=1{,}000$ and
$0.710$ at $n=10{,}042$. Those earlier numbers are therefore not comparable to these and
are not reported as such.

\begin{table}[t]
\centering\small
\caption{Full HellaSwag validation panel, 10{,}042 items, label-span likelihood.
$n=10{,}042$ for every row. Rows marked ($\dagger$) are the 0.4B controlled adaptation
matrix of Section~\ref{m:direction}, averaged over three training seeds, which are
\emph{not} decoding seeds. For these rows brackets give the training-seed range; other rows show recorded 95\% document-bootstrap intervals.}
\begin{tabular}{@{}llrr@{}}
\toprule
ID & Model & Accuracy & Interval\\
\midrule
F2 & BiRWKV diffusion, step 14{,}000 & 0.710 & $[0.701,0.719]$\\
--- & BiRWKV diffusion, step 13{,}000 & 0.715 & $[0.707,0.724]$\\
--- & Looped variants, steps 4{,}750--20{,}500 & 0.711--0.714 & ---\\
\midrule
R0 & RWKV-7 World3 2.9B (released causal) & \textbf{0.743} & $[0.734,0.751]$\\
--- & RWKV-7 World3 1.5B (released causal) & 0.686 & $[0.676,0.695]$\\
--- & RWKV-7 0.4B (released causal) & \textbf{0.534} & $[0.524,0.544]$\\
\midrule
$\dagger$ & 0.4B forward-only, no loop & 0.519 & $[0.517,0.522]$\\
$\dagger$ & 0.4B forward-only, with loop & 0.520 & $[0.518,0.521]$\\
$\dagger$ & 0.4B bidirectional, no loop & 0.511 & $[0.507,0.516]$\\
$\dagger$ & 0.4B bidirectional, with loop & 0.510 & $[0.507,0.515]$\\
\bottomrule
\end{tabular}
\label{tab:hellaswag}
\end{table}

The released causal comparators now carry this section's two substantive results. First,
at 2.9B the causal RWKV-7 reaches $0.743$ against the diffusion model's $0.710$: a
\textbf{3.3-point} gap in favour of the causal model, in the same direction as the 7.7-point MMLU gap. Second, and
more informative because it is controlled, the released \emph{causal} 0.4B checkpoint is
the initialization of the 0.4B adapted matrix, so the two sets compare with the
initialization held fixed. The released base scores $0.534$; the forward-only adapted arms
average $0.519$ and the bidirectional adapted arms $0.511$. The forward-only adapted arms are about 1.5 points below their initialization on HellaSwag, and the bidirectional arms about one point below the forward-only arms. The comparison to the initialization combines continued training, its data and objective; without equally budgeted causal continuation it does not isolate an objective effect.
These descriptive differences favour the released base on this panel; repeated-execution uncertainty for these arms has not been measured.

The 2.9B diffusion arms cluster within $0.5$ points of one another
($0.710$--$0.715$) across a fourfold range of training steps, which is the same
step-insensitivity the MMLU panel showed.

\subsection{Additional basic-ability panels}
\label{m:basic}
\begin{table}[t]
\centering
\small
\begin{tabular}{lrrr}
\toprule
Model & RACE & OpenBookQA & MMLU-Redux \\
Unique items & 4887 & 500 & 5700 \\
\midrule
F2, step 14000 & 46.90 & 35.80 & 38.67 \\
Loop, step 9500 & 46.61 & 36.60 & --- \\
Loop, step 20500 & 46.02 & 37.00 & --- \\
Released RWKV-7, 0.4B & 40.58 & 34.80 & 33.95 \\
Released RWKV-7, 1.5B & 45.98 & 39.40 & 40.21 \\
Released RWKV-7, 2.9B & 47.64 & 41.20 & 42.74 \\
\bottomrule
\end{tabular}
\caption{Basic-ability option-scoring accuracy (\%). RACE and OpenBookQA use validation splits; MMLU-Redux uses the test split. Each reported cell has exact unique item coverage at evaluation seed 42. Adapted checkpoints use \texttt{fwdce}; released checkpoints use \texttt{raw}. Dashes denote unavailable or unverified complete panels. These are descriptive checkpoint comparisons, with no claim of statistical significance or matched training budget.}
\label{tab:basic-ability-panels}
\end{table}


Table~\ref{tab:basic-ability-panels} reports complete RACE and OpenBookQA validation panels and available MMLU-Redux test panels. The collector checks each scoring arm against the exact unique document IDs in its task tree, rejects duplicate or incomplete coverage, and recomputes accuracy from the per-document records. Adapted models use the causal \texttt{fwdce} path and released models the equivalent \texttt{raw} path. These option-span likelihood panels are distinct from the one-token MMLU label protocol and are not pooled into a post-hoc aggregate.

F2 reaches 46.90\% on RACE and 35.80\% on OpenBookQA, compared with 47.64\% and 41.20\% for released RWKV-7 2.9B. The respective descriptive differences are $-0.74$ and $-5.40$ percentage points. F2 exceeds the 0.4B release on both panels and exceeds the 1.5B release on RACE. Differences in pretraining and adaptation remain confounded. F2 scores 38.67\% on MMLU-Redux versus 42.74\% for the 2.9B release. MMLU-Redux is a separate benchmark and its scores are not substitutes for full MMLU. Missing cells remain unreported until their complete records pass the same checks. Source hashes and coverage checks are recorded in \texttt{results/basic\_ability\_audit.json}; source hashes do not retroactively fill the upstream harness fields marked \texttt{unpinned}.

\begin{table}[t]
\centering
\small
\begin{tabular}{lrrrr}
\toprule
Dataset & $n$ & $\Delta$ (pp; 95\% CI) & F2 only & RWKV only \\
\midrule
RACE & 4887 & -0.74 [-1.74, +0.25] & 278 & 314 \\
OpenBookQA & 500 & -5.40 [-9.00, -2.00] & 28 & 55 \\
MMLU-Redux & 5700 & -4.07 [-5.04, -3.09] & 288 & 520 \\
\bottomrule
\end{tabular}
\caption{F2 step 14000 minus released RWKV-7 2.9B: accuracy differences in percentage points (pp), with 95\% paired-document percentile bootstrap intervals (10000 resamples; seed 20270921). The final columns count documents answered correctly only by the indicated model. Pairing uses exact document IDs at evaluation seed 42. Intervals describe item-sampling uncertainty for fixed checkpoints; pretraining and training budgets are unmatched, so differences are not causal effects.}
\label{tab:basic-ability-contrasts}
\end{table}

The paired intervals in Table~\ref{tab:basic-ability-contrasts} resample matched documents, with 10,000 replicates and a fixed seed. RACE is unresolved at this precision; the OpenBookQA and Redux differences favor the released comparator. These descriptive intervals are not multiplicity-adjusted confirmation tests and do not include between-training or between-execution uncertainty.

\subsection{Forward-only versus bidirectional, under matched training}
\label{m:direction}

Section~\ref{m:mechanism} lists "forward-only vs bidirectional" as a within-checkpoint
inference intervention: the same weights evaluated with the reverse stream removed. That
intervention isolates direction at inference. A different and stronger question is whether
\emph{training} with a bidirectional mixer helps, and the 0.4B controlled matrix answers
it, because all four arms share one initialization, one corpus and an equal token budget;
only the direction and the presence of a tied loop differ, at three training seeds each.

\begin{table}[t]
\centering\small
\caption{Controlled 0.4B arm matrix: direction $\times$ tied loop, from a common released
initialization, equal token budget, three training seeds per cell. Means over seeds;
brackets give the range across the three seeds. These are training-seed replicates of one
architecture family, not decoding repetitions of one checkpoint.}
\begin{tabular}{@{}llrr@{}}
\toprule
Direction & Loop & MMLU (14{,}042) & HellaSwag (10{,}042)\\
\midrule
\emph{released causal init} & --- & 0.2604 & 0.5338\\
\midrule
Forward-only & absent & 0.2686 $[0.2683,0.2691]$ & 0.5192 $[0.5166,0.5216]$\\
Forward-only & tied & 0.2689 $[0.2683,0.2695]$ & 0.5195 $[0.5178,0.5214]$\\
Bidirectional & absent & 0.2584 $[0.2531,0.2655]$ & 0.5108 $[0.5073,0.5156]$\\
Bidirectional & tied & 0.2583 $[0.2537,0.2636]$ & 0.5103 $[0.5071,0.5154]$\\
\bottomrule
\end{tabular}
\label{tab:direction}
\end{table}

The released causal initialization scores higher on HellaSwag than every adapted run. On MMLU, the forward-only adapted arms exceed the initialization, whereas the bidirectional means are below it. These are comparisons to an untouched initialization, not an equal-budget causal-continuation control, so they do not isolate the effect of changing the objective.

Within the adapted matrix, forward-only arms average 1.04 percentage points higher on MMLU and 0.88 points higher on HellaSwag than bidirectional arms. Tied-loop differences are small at this training budget: approximately $+0.000214$ and $-0.000047$ on MMLU, and $+0.0004$ and $-0.0005$ on HellaSwag, for forward and bidirectional models respectively. These observations concern the evaluated causal likelihood statistic and these checkpoints. They neither establish a benefit of bidirectionality nor determine its effect on long-context denoising. The three training seeds support descriptive within-family comparisons; they do not establish a general null effect for loops or a universal direction penalty.

\subsection{Execution variance bounds the resolvable contrast}
\label{m:variance}

Three of the readings above were scored twice, by accident of two campaign lanes, on the
same checkpoint and the same benchmark but on differently shaped pods: eight single-GPU
pods, and one eight-GPU pod. Comparing their per-item records gives a reproducibility
measurement of execution sensitivity. Changed item verdicts need not produce an equal change in aggregate accuracy.

\begin{table}[t]
\centering\small
\caption{The same checkpoint and benchmark scored under two pod placements. Per-item
comparison of verdicts and label likelihoods over the full 14{,}042-item MMLU panel.}
\begin{tabular}{@{}lrrr@{}}
\toprule
Checkpoint & Items changing verdict & Max $|\Delta \mathrm{NLL}|$ & Mean $|\Delta \mathrm{NLL}|$\\
\midrule
F2 (step 14{,}000) & 28 / 14{,}042 (0.20\%) & 0.104 & 0.0028\\
Looped step 12{,}500 & 105 / 14{,}042 (0.75\%) & 0.167 & 0.0080\\
Looped step 9{,}500 & 105 / 14{,}042 (0.75\%) & 0.216 & 0.0125\\
\bottomrule
\end{tabular}
\label{tab:variance}
\end{table}

The disagreement rates quantify item-level execution sensitivity, not the change in
aggregate accuracy: changes from correct to incorrect and from incorrect to correct can
cancel. Thus $0.75$ percentage points is an upper bound on the absolute accuracy change
for the two looped comparisons, not an estimated standard deviation or a universal
resolution floor. The displayed single-execution intervals do not include execution
variation, and a paired comparison with repeated executions is needed to assess a small
contrast. These three MMLU comparisons do not establish an execution threshold for
HellaSwag or for other checkpoints.

The mean absolute likelihood changes range from $3\times10^{-3}$ to
$1.3\times10^{-2}$. Their cause is unresolved; their magnitude alone cannot rule out
accumulated reduced-precision or kernel effects. The placements differ in pod width
(8$\times$1 versus 1$\times$8 GPUs), CPU count (20 versus 192, with no thread limit
set), and GPU/kernel configuration. Equal disagreement counts for two checkpoints do
not identify a shared cause without comparing the affected item sets and controls.

\subsection{Physical inference: memory and latency at 4K--64K}
\label{m:systems}

An exploratory single-device probe measures full-context forward passes, processed-input
throughput, and peak allocated and reserved device memory. It does not execute a
complete generation request. The saved BiRWKV rows contain no checkpoint digest or
resolved weight path; the adapter registry's default has 3.610B stored elements rather
than F2's 4.092B. Consequently, these measurements are attributed to the recorded
\texttt{birwkv-2.9b} probe label, not to the qualified F2 payload.

\begin{table}[t]
\centering\small
\caption{Exploratory single-device full-context forward probe, batch one, native
precision. Wall time is the sum of five timed calls after warm-up; throughput is
processed input tokens per second ($5L$/wall time), not generated tokens per second.
The NFE column is configured adapter metadata, not an observed sampler call count.
The BiRWKV checkpoint identity is unresolved. Two comparators failed to load.}
\begin{tabular}{@{}lrrrrr@{}}
\toprule
Model & NFE & Context & 5 calls (ms) & Input tok/s & Reserved (GiB)\\
\midrule
birwkv-2.9b (masked diffusion) & 100 & 4{,}096 & 1{,}205 & 16{,}990 & 9.43\\
 & & 16{,}384 & 4{,}624 & 17{,}716 & 13.46\\
 & & 32{,}768 & 9{,}199 & 17{,}810 & 18.86\\
 & & 65{,}536 & 18{,}187 & 18{,}017 & 29.58\\
\midrule
rwkv7-goose-world3-2.9b (AR) & 1 & 4{,}096 & 776 & 26{,}405 & 13.74\\
 & & 16{,}384 & 2{,}929 & 27{,}969 & 21.36\\
 & & 32{,}768 & 5{,}839 & 28{,}061 & 31.53\\
 & & 65{,}536 & 11{,}507 & 28{,}477 & 51.80\\
\midrule
llama-3.2-3b (AR transformer) & 1 & 4{,}096 & 313 & 65{,}331 & 7.81\\
 & & 16{,}384 & 1{,}544 & 53{,}068 & 13.40\\
 & & 32{,}768 & 3{,}763 & 43{,}543 & 20.73\\
 & & 65{,}536 & 10{,}242 & 31{,}995 & 35.42\\
\midrule
qwen2.5-3b (AR transformer) & 1 & 4{,}096 & 299 & 68{,}445 & 7.57\\
 & & 16{,}384 & 1{,}394 & 58{,}753 & 12.76\\
 & & 32{,}768 & 3{,}375 & 48{,}539 & 19.49\\
 & & 65{,}536 & 9{,}321 & 35{,}156 & 33.44\\
\midrule
llada-8b-base & --- & all & \multicolumn{3}{c}{load failed at every context}\\
mamba2-2.7b & --- & all & \multicolumn{3}{c}{load failed at every context}\\
\bottomrule
\end{tabular}
\label{tab:systems}
\end{table}

Over this measured range, processed-input throughput is approximately context-flat
for the recurrent mixers: the two RWKV rows rise by $6.0\%$ and $7.9\%$ between 4K
and 64K. Transformer throughput falls by factors of $2.04$ and $1.95$. This is a
forward-pass scaling observation, not a comparison of cached autoregressive generation
against iterative denoising.

Peak reserved memory grows from $9.43$ to $29.58$\,GiB for the BiRWKV probe, from
$13.74$ to $51.80$\,GiB for causal RWKV, and from $7.81$ to $35.42$\,GiB for Llama.
Fixed weight storage contributes to these totals, so growth by less than the $16\times$
length increase does not establish a sublinear asymptotic memory law. The BiRWKV 16K
row reserves $13.46$\,GiB on the tested device; this is not a physical 16-GiB-device
qualification and does not establish an F2 operating point.

Dividing processed-input throughput by the configured NFE cannot recover generated-token
throughput: the number of output tokens, actual sampler calls, changing canvases, and
causal prefill/decode work are not measured here. The $1.5\times$ goodput target requires
completed requests at a frozen quality threshold and deadline under concurrency, with
at least $1{,}000$ completions per cell. That target remains unmeasured. A fixed-overhead
estimate from two differently sourced canvas timings would likewise require matched
checkpoint, hardware, and execution settings before it could support an optimization
claim.

\subsection{Stress-test floors}
\label{m:floors}

The protocol requires free-generation floor effects to be reported rather than omitted.
F2 itself reads GSM8K at $0.61\%$ ($n=1{,}319$) across decoder variants, and earlier
runs of this lineage read GSM8K at or below $1.8\%$ (best over arms, $n=2{,}638$ and
$n=1{,}319$) and HumanEval pass@1 at or below $0.6\%$ ($n=328$). All are at floor; for
contrast the released LLaDA-8B attention-diffusion reference reads $61.5\%$, but at
$n=200$ and under a different decoder, so it is not the same statistic. These
are reported as context for the diffusion lineage, not as panel results under this
protocol: the item sets, the parser and the sampling schedule differ from the frozen
ones, and the executable-code suites the protocol names are the EvalPlus variants, which
are not scored here.

\subsection{Completed long-context confirmation}
\label{m:confirmation}

The locked 16K/32K/64K comparison is complete: 264 of 324 declared cells contain 200 paired examples each, giving 52,800 outcomes per model and 105,600 validated outcomes in total. The 60 generator refusals comprise six canvas overflows and 54 infeasible load-one cells. They remain explicit exclusions, not zero scores. Collection checks canonical instance IDs, prompt hashes, parsed answers against canonical gold, source provenance, and disjoint assigned coverage across execution partitions. Additional receipts outside a partition's assigned cells are counted but excluded rather than merged twice.

F2 uses the development-selected eight-stage matched-Bernoulli decoder at temperature $0.7$ and inference seed 17; R0 uses frozen greedy decoding. No confirmation outcome changes these policies. F2 answers 15 examples correctly and R0 answers 14; these pooled counts are distinct from the primary macro. Equal weighting of lengths, families, and supported cells within each family--length panel gives $0.030864\%$ for F2 and $0.028807\%$ for R0, a difference of $+0.002058$ percentage points. The descriptive paired 95\% interval for this difference is $[-0.020576,+0.024691]$ percentage points; the marginal intervals are $[0.016461,0.047325]\%$ and $[0.014403,0.045267]\%$, respectively. All recall and overwrite outcomes are incorrect. The only successes are in dataflow: at 32K, F2 has zero and R0 two; at 64K, F2 has 15 and R0 12. No model succeeds at 16K.

The intervals use 10,000 paired within-cell multinomial bootstrap draws with seed 20270921, preserving the primary macro's equal panel and cell weights. They condition on the generated cells, fixed checkpoints, and executed decoding seed, without representing unsupported-cell, training, or execution uncertainty. Zero-width intervals in all-zero panels do not establish equivalence. These results provide no evidence of long-context superiority or of the prespecified five-point practical gain. Together with the competence controls in Section~\ref{m:competence}, the near-zero endpoint does not isolate retention failure or a benefit from iterative refinement.

The validated aggregate and cell-level audit are recorded in \path{results/e3_confirmation_summary.json}. The portable export in \path{results/e3_receipt_export/manifest.json} records accepted per-item records, exact receipt hashes, source provenance, and explicit unsupported-cell records; its manifest tracks artifact counts and checksums separately from the aggregate.

\begin{table}[t]
\centering\small
\begin{tabular}{@{}llrrrr@{}}
\toprule
Model & Length & $n$ & NFE mean/med./p95 & Seconds med./p95 & Peak GiB alloc./res. \\
\midrule
F2 & 16K & 16,800 & 2.16/2.0/3.0 & 1.237/1.856 & 11.01/13.91 \\
F2 & 32K & 18,000 & 2.14/2.0/3.0 & 2.459/3.690 & 14.37/19.77 \\
F2 & 64K & 18,000 & 2.14/2.0/3.0 & 4.870/7.310 & 21.09/30.03 \\
R0 & 16K & 16,800 & 2.40/3.0/3.0 & 1.013/1.016 & 8.74/11.46 \\
R0 & 32K & 18,000 & 2.40/3.0/3.0 & 2.011/2.021 & 11.71/16.93 \\
R0 & 64K & 18,000 & 2.40/3.0/3.0 & 3.981/4.000 & 17.65/27.34 \\
\bottomrule
\end{tabular}
\caption{Descriptive per-request resources on completed supported confirmation cells. NFE counts actual model calls. Medians and p95 use inclusive linear interpolation; GiB values are maxima of request-level CUDA allocated/reserved bytes. F2 cache-replay qualification is outside item timing, whereas R0 first requests may include JIT compilation. R0 recomputes the full prefix with no KV cache. These measurements do not establish an optimized autoregressive speed comparison, speedup, or goodput.}
\label{tab:e3-confirmation-resources}
\end{table}

\begin{table}[t]
\centering\small
\begin{tabular}{@{}llrr@{}}
\toprule
Length & Family & F2 parseable/$n$ & R0 parseable/$n$ \\
\midrule
16K & Recall & 0/6,800 & 0/6,800 \\
16K & Dataflow & 18/5,000 & 26/5,000 \\
16K & Overwrite & 0/5,000 & 0/5,000 \\
32K & Recall & 0/7,200 & 0/7,200 \\
32K & Dataflow & 109/5,400 & 130/5,400 \\
32K & Overwrite & 0/5,400 & 0/5,400 \\
64K & Recall & 0/7,200 & 0/7,200 \\
64K & Dataflow & 1,385/5,400 & 632/5,400 \\
64K & Overwrite & 0/5,400 & 0/5,400 \\
\bottomrule
\end{tabular}
\caption{Task-grammar parsing on completed supported confirmation cells. Parseability is a diagnostic of output form and does not imply correctness. Unsupported cells are excluded, not zero-filled.}
\label{tab:e3-confirmation-parsing}
\end{table}


\subsection{Development-only decoder selection}
\label{m:devselection}
The disjoint development set contains 72 examples: 24 each of recall, overwrite, and dataflow at 16K tokens, position 50\%, binding load 8, and similar distractors. Its data seeds are 201--203; confirmation uses 101--105 and a separate generator split. All 16 F2 configurations score 0/72 exact matches. The frozen tie rule therefore selects matched-Bernoulli revelation, eight nominal stages, and temperature $0.7$. This is a tied operating point, not evidence that it is superior. The released causal RWKV-7 comparator also scores 0/72 under greedy decoding. These outcomes motivate separately reported task-competence controls; they do not authorize tuning on confirmation data.

\begin{table}[h]
\centering\small
\begin{tabular}{@{}lrrrr@{}}
\toprule
Reveal policy & Stages & Temperatures & Correct / examples & Mean actual NFE\\
\midrule
Bernoulli & 8 & $0.7,1.0$ & $0/72$ each & 2.111\\
Bernoulli & 16 & $0.7,1.0$ & $0/72$ each & 2.208\\
Bernoulli & 32 & $0.7,1.0$ & $0/72$ each & 2.292\\
Bernoulli & 64 & $0.7,1.0$ & $0/72$ each & 2.306\\
Confidence & 8 & $0.7,1.0$ & $0/72$ each & 2.333\\
Confidence & 16 & $0.7,1.0$ & $0/72$ each & 2.333\\
Confidence & 32 & $0.7,1.0$ & $0/72$ each & 2.333\\
Confidence & 64 & $0.7,1.0$ & $0/72$ each & 2.333\\
\bottomrule
\end{tabular}
\caption{Development sweep. Each row combines two separately executed temperatures with identical aggregate outcomes and NFE. No unchanged canvas requires another network call after its logits are cached.}
\label{tab:devselection}
\end{table}

Caching is limited to target logits for a byte-identical current canvas. Any committed token invalidates that cache; all sampling draws and schedule steps are preserved. Sixteen qualification checks (eight ranks, two sampler families) compare cached and uncached trajectories and require identical output IDs and per-stage commit counts. Repeated uncached calls on the same canvas also produce bitwise-equal target logits in these checks. All checks pass. Nominal reverse stages and actual neural evaluations are therefore reported separately: the selected configuration uses 152 calls across 72 requests, or 2.111 calls per request. This optimization does not constitute persistent recurrent-history reuse and does not establish a deployment speedup.

The output supports are explicitly different: F2 excludes its mask and padding IDs (65535 and 0) from sampling but retains other head entries, while R0 restricts greedy output to the 65,529 qualified native byte-token IDs. No chat template is applied. Both receive the gold span length as the generation budget; an extra leading token can therefore prevent exact completion. These frozen policies are properties of the evaluated interfaces, not a matched output-support ablation or recovery of the historical grouped decoder.


\subsection{Task-format competence controls}
\label{m:competence}
After decoder freezing, we rerun the same 72 development tasks with only filler removed, with evidence replaced by neutral filler, and with the same logical tasks rerendered at 4K and 8K. These are post-development diagnostics, not independent confirmation. Gold target length remains supplied; the decoder and answer parser are unchanged. Table~\ref{tab:e3-controls} reports all 576 outcomes.

Both models score zero and produce no parseable answers in all four controls. In particular, removing all irrelevant filler leaves only 175--180 tokens but does not restore successful task execution. Consequently, a long-context floor under this serialization cannot isolate memory retention or a benefit from iterative refinement. The observed zero is specific to this prompt, supplied-length generation interface, and fixed checkpoints; it is not a general claim that the models lack recall or dataflow capability. The 4K/8K diagnostics cannot be substituted for full-grid degradation contrasts.

\begin{table}[t]
\centering
\small
\begin{tabular}{llrrrr}
\toprule
Control & Family & F2 correct/$n$ & R0 correct/$n$ & F2 parse (\%) & R0 parse (\%) \\
\midrule
Evidence only & Recall & 0/24 & 0/24 & 0.0 & 0.0 \\
Evidence only & Overwrite & 0/24 & 0/24 & 0.0 & 0.0 \\
Evidence only & Dataflow & 0/24 & 0/24 & 0.0 & 0.0 \\
 & Family macro (\%) & 0.00 & 0.00 & 0.0 & 0.0 \\
\midrule
No evidence (16K) & Recall & 0/24 & 0/24 & 0.0 & 0.0 \\
No evidence (16K) & Overwrite & 0/24 & 0/24 & 0.0 & 0.0 \\
No evidence (16K) & Dataflow & 0/24 & 0/24 & 0.0 & 0.0 \\
 & Family macro (\%) & 0.00 & 0.00 & 0.0 & 0.0 \\
\midrule
Paired 4K & Recall & 0/24 & 0/24 & 0.0 & 0.0 \\
Paired 4K & Overwrite & 0/24 & 0/24 & 0.0 & 0.0 \\
Paired 4K & Dataflow & 0/24 & 0/24 & 0.0 & 0.0 \\
 & Family macro (\%) & 0.00 & 0.00 & 0.0 & 0.0 \\
\midrule
Paired 8K & Recall & 0/24 & 0/24 & 0.0 & 0.0 \\
Paired 8K & Overwrite & 0/24 & 0/24 & 0.0 & 0.0 \\
Paired 8K & Dataflow & 0/24 & 0/24 & 0.0 & 0.0 \\
 & Family macro (\%) & 0.00 & 0.00 & 0.0 & 0.0 \\
\bottomrule
\end{tabular}
\caption{Post-development competence diagnostics using the same 72 development tasks (24 per family) for both models and all controls. R0 is released RWKV-7 2.9B; F2 uses its frozen 8-step, temperature-0.7 matched-Bernoulli decoder. Evidence-only removes marked filler; no-evidence replaces the evidence payload with neutral filler at 16K. Paired 4K/8K anchors rerender the identical logical task without redrawing its bindings. Exact-match counts, family-macro accuracy and parsing rates retain the supplied target length. These reused-development diagnostics are not independent confirmation and do not estimate full-grid $\Gamma$.}
\label{tab:e3-controls}
\end{table}


\subsection{Exploratory generation canaries}
\label{m:canary}

To inspect the development-task floor, we subsequently ran eight fixed prompts per model spanning capital completion, arithmetic, lookup instructions, one-shot lookup, symbolic serialization, and chat-style variants. Each run supplies a 32-token masked suffix, replacing the short gold-length answer budget used in the symbolic task evaluations. F2 retains the 8-step, temperature-0.7 matched-Bernoulli decoder; R0 retains greedy decoding. A separate first-position diagnostic adds one forward call to each run. These 16 records are exploratory observations, with no post-hoc accuracy threshold or aggregate success score.

Table~\ref{tab:e3-canary} shows leading excerpts from three paired prompts. Both capital-chat continuations mention Paris. The R0 chat-lookup continuation includes the stored reference string, whereas the corresponding F2 continuation does not provide that value in its leading excerpt. The raw symbolic serialization produces the displayed fragmented or repeated continuations. These examples show that the earlier floor is compatible with nontrivial continuations under a different interface; they do not establish a prompt-only explanation, general task competence, or an improvement in long-context retention. Prompt wording and output budget change together, and these hand-inspected outputs neither replace nor tune the locked confirmation endpoint.

The audit checks all eight prompt pairs, regenerated masked-canvas hashes, 32-token output budgets, output-ID decoding, rank provenance, and accessible source and weight hashes. It retains every output and diagnostic in \texttt{results/e3\_canary\_audit.json}. Runtime-library hashes from the job are distinguished from checks of files accessible on the audit host: matching source bytes corroborates content but does not recreate the historical container environment.

\begin{table}[t]
\centering\small
\begin{tabular}{@{}p{0.18\linewidth}p{0.38\linewidth}p{0.38\linewidth}@{}}
\toprule
Prompt & F2 leading excerpt & R0 leading excerpt \\
\midrule
Capital, chat & \texttt{It is Paris. Paris city name} & \texttt{Paris.} \\
Lookup, chat & \texttt{the, the.  ----> select the3 to the3.} & \texttt{The value stored for key093 is v1911.} \\
Symbolic raw & \texttt{a:\textbackslash{}nc:\textbackslash{}n2\textbackslash{}nd:\textbackslash{}n1} & \texttt{dvk093\textbackslash{}ndvk093\textbackslash{}ndvk093} \\
\bottomrule
\end{tabular}
\caption{Illustrative leading excerpts from exploratory 32-token canaries; leading whitespace is omitted and literal \texttt{\textbackslash{}n} denotes a newline. Excerpts are not a scoring rule. The paired models receive identical prompts and budgets; all eight prompt pairs, full output IDs/text and source hashes are retained in \texttt{results/e3\_canary\_audit.json}.}
\label{tab:e3-canary}
\end{table}


\subsection{What is not measured}
\label{m:notmeasured}

The following distinguish unexecuted parts of the broader design from the measured subset. They are not lower bounds or inferred results.

\begin{itemize}[leftmargin=1.2em,itemsep=2pt,topsep=2pt]
\item \textbf{Full-grid paired $4$K/$8$K anchors and $\Gamma_b(L)$}: not executed. The smaller post-development competence diagnostics in Section~\ref{m:competence} retain the same 72 development tasks and cannot replace these full-grid contrasts.
\item \textbf{External suites} (RULER, LongBench, BABILong, NoLiMa): not run. RULER's
pinned grid is 4K/8K/16K and is generated locally; LongBench and BABILong would each need
their own access-matched serialization for causal baselines, which the protocol requires
and which does not yet exist.
\item \textbf{L0 (GLA) and L1 (DeltaNet) long-context comparisons}: neither is executed. GLA did, however, complete a separate causal MMLU test panel: 3,232 correct out of 14,042 unique items (23.02\%). The eight raw-scoring shards agree exactly with their merged records and the tokenizer-specific test tree; the resolved configuration declares \texttt{GLAForCausalLM}. Thus earlier architecture-registration and tokenizer-loading failures are historical blockers, not evidence that no hardware load succeeded. The model snapshot is \path{46b15820a4df269e99aed9d709e017677c15d24b}; upstream harness provenance fields remain unpinned. The audit in \texttt{results/additional\_baseline\_audit.json} records coverage, source hashes, and remaining limitations. This capability observation does not supply a long-context GLA result. DeltaNet has no complete campaign capability or long-context panel; the absent payload is not replaced by a different model.
\item \textbf{M0 (Mamba-2 2.7B)}: absent. The only Mamba reading in the campaign is a
Mamba-1 2.8B checkpoint, whose configuration declares the Mamba-1 architecture class;
the Mamba-2 2.7B checkpoint named in Table~\ref{m:arms} exists but has no loader in
this harness and no reading. A Mamba-1 result is not a Mamba-2 result and we do not
substitute one for the other.
\item \textbf{D0 (LLaDA-8B-Base)}: the checkpoint is present and has a separate loader
outside this harness; no full-panel MMLU reading for it exists under this protocol, and
it failed to load in the systems probe. Its only MMLU observation is the $n=200$ item of
the historical panel below, which the protocol's own text forbids reading as full MMLU.
\item \textbf{Incomplete auxiliary cells}: complete RACE and OpenBookQA panels are now reported in Section~\ref{m:basic}. Remaining missing model--dataset combinations are explicitly marked in Table~\ref{tab:basic-ability-panels}. Historical panels at different item counts or prompt constructions are not substituted for these cells.
\item \textbf{The historical card-style capability panel} (MMLU at $n=200$, OpenBookQA
and RACE at $n=500$, LLaDA-8B and the released RWKV-7 among its rows): measured, but
under a different protocol and a different item count, and the protocol's own text
forbids reading a 200-question observation as full MMLU. It is context, not a result of
this study, and the $n=200$ MMLU score is additionally extractor-dependent --- a
first-token match and a substring match on the same generations give $0.095$ and $0.230$
respectively, so the number is not stable enough to report even as context.
\item \textbf{Native infilling} (Section~\ref{m:quality}): not run. It requires an
information-access-matched serialization for causal baselines, and reporting it with a
prefix-only baseline would be the access mismatch the protocol forbids.
\item \textbf{Goodput, concurrency and the low-memory variants}: the single-device
memory--latency frontier is measured (Section~\ref{m:systems}), but the protocol's
$1.5\times$ goodput target needs concurrency, a frozen deadline and $\ge1{,}000$
completions per cell, and none of those were run. Quantization, selective heads and
offload are named modifications that need parity and quality checks and were not
attempted. Two of the six models (LLaDA-8B and Mamba-2 2.7B) failed to load in the probe
and have no systems row at all.
\item \textbf{State-only retention, cross-backbone transfer, and causal architectural
attribution}: not claimed, and no experiment here supports them.
\end{itemize}

